\chapter*{ВВЕДЕНИЕ}
\addcontentsline{toc}{chapter}{Введение}

Цифровые инструменты стали частью подготовки и сопровождения туристических
маршрутов: пользователь заранее изучает карту, фотографии, отзывы, сложность и
условия прохождения. По данным Всемирной туристской организации, число
международных туристических поездок к 2024~году превысило 1,3~млрд в год, а
значительная доля путешественников использует цифровые сервисы на этапе
планирования и во время поездки~\cite{unwto2024}. В Российской Федерации
туристская деятельность регулируется федеральным законодательством~\cite{tourismLaw132},
а стратегия развития туризма до 2035~года рассматривает цифровизацию как один
из инструментов повышения доступности и качества туристских услуг~\cite{tourismStrategy2035}.

В рамках данной работы рассматривается не весь рынок туристических сервисов, а
конкретный пользовательский сценарий: автор маршрута должен подготовить
интерактивный материал о прогулке, экскурсии или походе. Такой материал включает
линию маршрута, точки, фотографии, пояснения, сведения о сложности и способ
публикации для потенциальных участников. Для гида, инструктора или небольшой
туристической организации это прикладная задача: маршрут нужно не только
сохранить для себя, но и представить клиенту как цельный цифровой объект.

Ключевая особенность такого объекта заключается в связи фотографий с координатами.
Фотография обзорной точки, достопримечательности, развилки или сложного участка
имеет практический смысл только тогда, когда пользователь видит её положение на
карте и связь с соседними точками маршрута. Поэтому в работе используется понятие
геопривязанной фотографии: фотографии, связанной с координатами точки маршрута
вручную пользователем или автоматически на основе GPS-данных из EXIF-метаданных.

Существующие решения закрывают эту задачу частично. Картографические инструменты,
например Google My Maps~\cite{googlemymaps}, позволяют создавать пользовательские
карты, но не ориентированы на полный цикл подготовки туристического маршрута:
работу с фототочками, публикацию, обсуждение, оценки и повторное использование.
Сервисы активного туризма, такие как Wikiloc~\cite{wikiloc}, Strava~\cite{strava},
Komoot~\cite{komoot} и AllTrails~\cite{alltrails}, в большей степени ориентированы
на запись, навигацию или каталог треков. Подробное сравнение аналогов по
18~критериям приведено в разделе~\ref{sec:competitors}.

Прикладная проблема состоит в том, что автору маршрута приходится использовать
несколько разрозненных инструментов: карту для линии маршрута, файловое хранилище
или мессенджер для фотографий, отдельный текстовый документ для описания и ещё
один канал для публикации ссылки клиентам. В результате теряется связь между
координатами, фотографиями и пояснениями, а повторное использование готового
маршрута становится трудоёмким.

Актуальность работы определяется следующими факторами:
\begin{itemize}
  \item ростом спроса на цифровые инструменты планирования и представления
    туристических маршрутов;
  \item потребностью гидов, инструкторов и организаторов походов в едином
    инструменте подготовки маршрута с фотографиями и пояснениями;
  \item недостаточной поддержкой сценария управляемой публикации авторского
    маршрута с фототочками, комментариями и оценками в существующих сервисах;
  \item необходимостью проверяемой реализации, в которой маршрут, фотографии,
    социальные функции и административное управление объединены в одной
    информационной системе.
\end{itemize}

Целью данной работы является разработка корпоративной информационной системы
картографии и маршрутов Guide Helper, предназначенной для картографического
описания, создания и публикации туристических маршрутов с геопривязанными
фотографиями.

Для достижения поставленной цели необходимо решить следующие задачи:
\begin{enumerate}
  \item проанализировать предметную область цифрового туризма и потребности
    целевой аудитории;
  \item выполнить обзор и сравнительный анализ существующих аналогов;
  \item сформулировать функциональные и нефункциональные требования к системе;
  \item спроектировать архитектуру системы и схему хранения данных;
  \item реализовать серверные сервисы для аутентификации, маршрутов и обработки
    фотографий;
  \item разработать клиентские интерфейсы для создания, просмотра и публикации
    маршрутов;
  \item реализовать обработку геопривязанных фотографий и расчёт характеристик
    маршрута;
  \item реализовать интерактивные функции: комментарии, оценки, закладки,
    уведомления и ИИ-ассистента;
  \item настроить инфраструктуру сборки, проверки и развёртывания системы;
  \item провести тестирование системы и оценить соответствие результата
    сформулированным требованиям.
\end{enumerate}

Объектом исследования является процесс создания, хранения и распространения
туристических маршрутов с геопространственными данными и мультимедийным
контентом.

Предметом исследования являются методы и средства разработки информационной
системы для работы с туристическими маршрутами, геопривязанными фотографиями и
пользовательским взаимодействием.

В работе применяются следующие методы: анализ предметной области и аналогов,
формирование требований, проектирование архитектуры программной системы,
моделирование данных, разработка серверной и клиентской частей, функциональное и
модульное тестирование, анализ производительности и проверка пользовательских
сценариев.

Практическая значимость работы заключается в создании работающей системы,
которая уменьшает число разрозненных действий при подготовке маршрута. Автор
маршрута может в одном интерфейсе создать маршрут, добавить точки, прикрепить к
ним фотографии и заметки, опубликовать ссылку и получить обратную связь. Турист
может открыть маршрут как интерактивную карту, увидеть фотографии в привязке к
местам, оценить сложность и принять решение о прохождении маршрута. Реализация
проверена по 37~функциональным требованиям, 115~модульным тестам и сценариям
функционального тестирования.

Работа состоит из введения, трёх глав, заключения, списка использованных
источников и трёх приложений: технического задания, программы и методики
испытаний, материалов сравнительной оценки аналогов.

В первой главе анализируются предметная область, целевая аудитория и аналоги.
На основе анализа формулируются функциональные и нефункциональные требования к
системе, обосновывается выбор технологического стека и описываются алгоритмы
расчёта расстояния, времени и сложности маршрута.

Во второй главе проектируются архитектура, схема базы данных и макеты основных
экранов. В этой же главе рассматривается реализация системы: сервис
аутентификации, сервис маршрутов, асинхронный обработчик фотографий,
ИИ-ассистент, клиентское веб-приложение, схема хранения данных и инфраструктура
развёртывания.

В третьей главе приводятся результаты модульного, функционального и
нагрузочного тестирования, проверяется выполнение требований и рассматривается
экономическое обоснование разработки.
